\section{Core Rules for Thiasos}

A Thiasos is a fragment of a Patron's attention made manifest. It is not a pet, not a servant, not a tool. It is a bond—a living, breathing, or ticking witness to the Runekeeper's covenant. Every Runic Warrior has a Thiasos, and every Thiasos has its own nature, its own needs, and its own memory.

\subsection{What a Thiasos Is}

A Thiasos is a companion bound to the Runekeeper through their shared covenant. It serves as:

\begin{itemize}
\item \textbf{A Witness:} The Thiasos sees and remembers every Rite the Runekeeper invokes. It counts as a witness for any Rite that requires one.
\item \textbf{A Helper:} The Thiasos may assist the Runekeeper in physical tasks, combat, or perception. It acts on the Runekeeper's turn unless otherwise noted.
\item \textbf{A Burden:} The Thiasos can be wounded, captured, or destroyed. Losing a Thiasos is a loss of witness and a mark of failure.
\end{itemize}

A Thiasos is not a separate character. It does not have its own agenda (unless the GM decides otherwise for narrative effect). It exists to serve the covenant—but it may have opinions.

\subsection{Thiasos Traits}

Every Thiasos has the following traits:

\subsubsection{Cap (Capability)}

A measure of the Thiasos's power and effectiveness, typically ranging from 2 to 5. Cap determines:

\begin{itemize}
\item How many dice the Thiasos adds when assisting (min(Cap, relevant skill)).
\item The difficulty of tasks the Thiasos can attempt independently.
\item The Thiasos's Harm threshold (typically Cap × 1.5, rounded up).
\end{itemize}

\subsubsection{Harm}

A Thiasos can take Harm. When its Harm track is full, it is Incapacitated. An Incapacitated Thiasos cannot act, cannot witness, and does not recover on its own. Healing a Thiasos requires:

\begin{itemize}
\item \textbf{Living Thiasos:} Rest, Medicine, or magical healing (same as a character).
\item \textbf{Construct Thiasos:} Repair (Craft test, DV 3-5, downtime).
\item \textbf{Spirit Thiasos:} A ritual at a threshold (Spirit + Lore, DV 4, one hour).
\item \textbf{Sentient Thiasos:} Negotiation. The Thiasos may demand a service or a memory before it accepts repair.
\end{itemize}

\subsubsection{Traits}

Special abilities unique to the Thiasos's type (warhorse, spirit-hound, sentient blade, etc.). Traits are described in the individual Thiasos entries.

\subsection{Assisting the Runekeeper}

A Thiasos may assist the Runekeeper on a roll by providing assistance dice. The number of dice equals min(Cap, relevant skill), to a maximum of +3 total from all sources (including other allies). The Runekeeper must describe how the Thiasos helps.

\paragraph{Examples:}
\begin{itemize}
\item A warhorse grants +2 dice to a Ride roll.
\item A spirit-hound grants +1 die to a Tracking roll.
\item A sentient blade grants +1 die to a Parry roll.
\end{itemize}

A Thiasos may also act independently once per scene, taking its own action on the Runekeeper's turn. Independent actions do not require the Runekeeper to spend their action commanding. The Thiasos's action uses its own Cap and traits.

\subsection{Witness and Testimony}

A Thiasos counts as a witness for any Rite that requires one. It sees, it hears, it remembers. In a hedge-court or religious tribunal, a Thiasos's testimony may be accepted—though its form of communication may be limited (growls, chirps, bell tones, written words on a slate).

A sentient Thiasos (one that can speak) can provide full testimony. Its word is considered as binding as a living witness, though some courts may question its impartiality.

\subsection{The Bond}

The Runekeeper and Thiasos are bound. This bond has several implications:

\paragraph{Shared Perception:} The Runekeeper may sense the Thiasos's location (within Far range) and its emotional state. The Thiasos may sense the Runekeeper's physical state (wounded, exhausted, dying).

\paragraph{Shared Burden:} With the Shared Burden talent, the Runekeeper may transfer Fatigue to the Thiasos. The Thiasos becomes Winded (-1 die to all actions) for the rest of the scene.

\paragraph{Shared Obligation:} When the Runekeeper marks Obligation, the Thiasos may also bear a visible mark. A warhorse develops a limp. A spirit-hound's glow dims. A sword's edge dulls. The mark fades when the Obligation is cleared.

\subsection{Losing a Thiasos}

A Thiasos can be killed, destroyed, or lost.

\begin{itemize}
\item If a \textbf{living Thiasos} dies, the Runekeeper marks 2 Obligation immediately and gains the Condition \textit{Bereaved} (-1 die to all rolls involving trust or loyalty) until they acquire a new Thiasos.
\item If a \textbf{construct or spirit Thiasos} is destroyed, the Runekeeper loses all Rites engraved on it (if it served as a Codex). They may repair it with significant downtime and a Craft test (DV 5).
\item If a \textbf{sentient Thiasos} is destroyed, the Runekeeper loses its counsel and its witness. A sentient Thiasos cannot be replaced—only mourned.
\end{itemize}

\subsection{Acquiring a New Thiasos}

Acquiring a new Thiasos is a downtime project. The Runekeeper must:

\begin{enumerate}
\item Find a suitable creature, construct, or object (GM discretion).
\item Perform a bonding ritual (one day, requires a threshold and a witness).
\item Mark 1 Obligation (the Patron's fee for a new ally).
\end{enumerate}

A new Thiasos starts with Cap 1 and no traits. It gains Cap and traits as the Runekeeper advances (see Living Thiasos and Sentient Thiasos chapters).

\subsection{Thiasos in Combat}

A Thiasos uses the standard combat rules with the following adjustments:

\begin{itemize}
\item \textbf{Initiative:} The Thiasos acts on the Runekeeper's turn. It may act before or after the Runekeeper.
\item \textbf{Movement:} The Thiasos moves independently. It does not need to stay Close to the Runekeeper, but if it moves beyond Far range, the bond is strained. The Runekeeper cannot sense the Thiasos's location, and the Thiasos cannot witness Rites.
\item \textbf{Defense:} The Thiasos can be targeted separately. It uses its own Body and Athletics (or equivalent) for defense.
\item \textbf{Incapacitation:} An Incapacitated Thiasos cannot act. It does not die immediately, but if left untreated (or unrepaired) for more than a day, it dies or is destroyed.
\end{itemize}

\subsection{Thiasos and Codex}

A sentient Thiasos may also serve as the Runekeeper's Codex. Rites can be engraved, etched, or woven into its surface. This is particularly thematic for sentient blades, speaking staves, and living masks.

If a Thiasos that serves as a Codex is destroyed, the Runekeeper loses all Rites engraved on it. They may retrieve the Rites from memory (requiring a Spirit + Lore test, DV 4, one day per Rite) or re-engrave them from a backup.

\subsection{Thiasos Advancement}

A Thiasos advances when the Runekeeper advances in Tier. At each new Tier, the Thiasos gains:

\begin{itemize}
\item +1 Cap (maximum 5).
\item One new trait (from its type list).
\item +2 Harm capacity.
\end{itemize}

Some Thiasos may gain unique traits based on the Runekeeper's Patron or the Thiasos's history. The GM and player should collaborate.

\subsection{Quick Reference}

\begin{tabular}{ll}
\textbf{Feature} & \textbf{Rule} \\
\hline
Assist Dice & min(Cap, relevant skill), max +3 total \\
Independent Action & Once per scene, on Runekeeper's turn \\
Witness & Counts for all Rites \\
Shared Burden & Transfer Fatigue (requires talent) \\
Loss & Mark 2 Obligation + Bereaved Condition \\
Acquisition & Downtime project, 1 Obligation, requires witness \\
Advancement & +1 Cap, +1 trait, +2 Harm per Tier \\
\end{tabular}

\subsection{Final Note}

A Thiasos is not a weapon. It is not a shield. It is a witness, a companion, a fragment of covenant made manifest. Treat it well. It remembers.

The thread held. The water carried. That is not a Rite. That is grace.
